%===============================================================================
% Template para artigo CBA/SBA - Adaptado do IFAC
% Artigo: Controle e Automação da Planta Piloto de Hidrogênio por
%         Reforma a Seco de Biogás – UFPR/LabMater
%
% ORDEM DE COMPILAÇÃO:
%   1. pdflatex main.tex
%   2. bibtex main
%   3. pdflatex main.tex
%   4. pdflatex main.tex
%
% CHAVES DO ifacconf.bib USADAS NESTE ARQUIVO:
%   oliveira2021, oliveira2022, escribano2025, lavoie2014, luyben2014,
%   mokashi2024, essien2019, thermodynamics2025, labmater2024,
%   dissertacao2024, escribano_ipt, balancing2023,
%   microwave2025, photothermal2024
%===============================================================================
\documentclass[a4paper]{ifacconf}

\usepackage{graphicx,amsmath,url}
\usepackage[round]{natbib}

\def\portugues{1}

\ifx\portugues\undefined
\def\portugues{0}
\fi

\if\portugues0
   \usepackage[english]{babel}
  \else
   \usepackage[spanish,brazil,english]{babel}
\fi

\usepackage[T1]{fontenc}
\usepackage[utf8]{inputenc}
\usepackage{ae}

\if\portugues1
 \newtheorem{teorema}[thm]{{\em Teorema}}{ }
 \newtheorem{lema}[thm]{{\em Lema}}{ }
 \newtheorem{corolario}[thm]{{\em Corolário}}{ }
 \newenvironment{prova}{{\bf Prova.}}{ }
\fi

\begin{document}

\if\portugues1
\selectlanguage{brazil}

\begin{frontmatter}

\title{Controle e Automação de Planta Piloto de Produção de
Hidrogênio por Reforma a Seco de Biogás no
LabMater/UFPR\thanksref{footnoteinfo}}

\thanks[footnoteinfo]{Reconhecimento do suporte financeiro deve vir
nesta nota de rodapé.}

\author[First]{Ricardo M. Israel}
\author[Second]{Dr. Gustavo H. C. Oliveira}
\author[Third]{Dr. Helton J. Alves}

\address[First]{Departamento de Engenharia Elétrica, Universidade
Federal do Paraná, Curitiba, Paraná, Brasil (e-mail: autor1@ufpr.br).}
\address[Second]{Departamento de Engenharia Elétrica, Universidade
Federal do Paraná, Curitiba, Paraná, Brasil (e-mail: autor2@ufpr.br).}
\address[Third]{Departamento de Engenharia Química, Universidade
Federal do Paraná, Curitiba, Paraná, Brasil (e-mail: autor3@ufpr.br).}

\selectlanguage{english}
\renewcommand{\abstractname}{{\bf Abstract:~}}

\begin{abstract}
This paper presents the design and experimental validation of a process
control architecture for a pilot-scale renewable hydrogen production plant
based on dry reforming of \emph{biogas} (BDR), developed at the Laboratory
of Materials and Renewable Energies (LabMater) of the Federal University of
Paraná (UFPR). In the studied process, raw biogas---a CH$_4$/CO$_2$ mixture
with variable composition and typical impurities---is converted into synthesis
gas (H$_2$ + CO) via highly endothermic reactions, which makes the reactor
strongly sensitive to heat supply, feed ratio and flow disturbances. These
characteristics, combined with the risk of coke deposition and catalyst
deactivation, impose strict operational constraints that require automatic
control of reactor temperature and CO$_2$/CH$_4$ feed ratio within a safe
operating window. The paper formulates the mathematical control problem for
the LabMater/UFPR biogas dry reforming pilot plant, proposes a control
structure tailored to this configuration, and validates it experimentally.
Results show that the implemented architecture maintains the reactor in the
target temperature range, ensures high CH$_4$ conversion and keeps the
CO$_2$/CH$_4$ ratio within limits that mitigate coke formation, illustrating
the specific challenges and benefits of controlling biogas, rather than pure
methane, in dry reforming.
\end{abstract}

\selectlanguage{brazil}
\renewcommand{\abstractname}{{\bf Resumo:~}}

\begin{abstract}
Este artigo apresenta o projeto e a validação experimental de uma arquitetura
de controle de processos para uma planta piloto de produção de hidrogênio
renovável baseada na reforma a seco de \emph{biogás}, desenvolvida no
Laboratório de Materiais e Energias Renováveis (LabMater) da UFPR. No processo
estudado, o biogás bruto---mistura CH$_4$/CO$_2$ de composição variável e
contendo impurezas típicas---é convertido em gás de síntese (H$_2$ + CO) por
meio de reações fortemente endotérmicas, o que torna o reator sensível a
perturbações no suprimento de calor, na razão CO$_2$/CH$_4$ e nas vazões de
alimentação. Essas características, somadas ao risco de deposição de coque e
desativação do catalisador, impõem restrições operacionais rígidas e demandam
controle automático preciso da temperatura do leito e da razão de alimentação
para manter a operação em uma janela segura. O trabalho formula o problema de
controle para a planta piloto de reforma a seco de biogás do LabMater/UFPR,
propõe uma estrutura de controle adequada a essa configuração e a valida
experimentalmente, evidenciando os desafios adicionais do uso de biogás, em
comparação ao metano puro, na operação e automação de processos de reforma a
seco.
\end{abstract}

\selectlanguage{english}

\begin{keyword}
Biogas dry reforming; green hydrogen production; process control;
pilot plant; LabMater/UFPR.

\vskip 1mm
\selectlanguage{brazil}
{\noindent\it Palavras-chaves:} Reforma a seco de biogás; reforma a seco do
metano; produção de hidrogênio verde; controle de processos; planta piloto;
LabMater/UFPR.
\end{keyword}

\selectlanguage{brazil}

\end{frontmatter}
\fi

%===============================================================================
%  CORPO DO ARTIGO
%===============================================================================

\section{Introdução}

A reforma a seco de biogás (BDR, do inglês \textit{biogas dry reforming})
consiste na conversão de uma mistura CH\(_4\)/CO\(_2\) proveniente de digestão
anaeróbia em gás de síntese (H\(_2\) + CO), utilizando o próprio CO\(_2\) do
biogás como agente oxidante do metano.\citep{lavoie2014} Ao substituir o uso
de metano puro por biogás, o processo passa a atuar diretamente sobre uma
corrente renovável e de menor valor agregado, contribuindo simultaneamente para
a valorização de resíduos orgânicos e para a mitigação de emissões de gases de
efeito estufa.\citep{lavoie2014} Nesse contexto, a reforma a seco de biogás
se destaca como rota promissora para a produção de hidrogênio renovável e gás
de síntese integrados a cadeias de energia de baixo carbono.

Do ponto de vista termodinâmico, a reação global de reforma a seco é fortemente
endotérmica, com entalpia padrão da ordem de 247~kJ·mol\(^{-1}\), de modo que
temperaturas elevadas e suprimento intenso de calor são necessários para
alcançar conversões significativas de CH\(_4\) e CO\(_2\).\citep{thermodynamics2025}
Estudos de equilíbrio mostram que, para misturas representativas de biogás, a
conversão de metano aumenta com a temperatura e que faixas típicas de operação
situam-se entre 650 e 900~\(^\circ\)C, enquanto abaixo de cerca de 640~\(^\circ\)C
a reação deixa de ser espontânea e a formação de carbono sólido torna-se
favorecida.\citep{thermodynamics2025,lavoie2014} Em comparação com a reforma a
seco de metano puro, o uso de biogás introduz complexidades adicionais, como a
variabilidade temporal da composição CH\(_4\)/CO\(_2\), a possível presença de
espécies inertes (N\(_2\), O\(_2\)) e de traços de impurezas (por exemplo,
H\(_2\)S), que afetam tanto o desempenho catalítico quanto o comportamento
dinâmico do processo.

Apesar da literatura extensa em desenvolvimento de catalisadores, estudos
cinéticos e análises termodinâmicas para reforma a seco, há relativamente
poucos trabalhos dedicados a controle e automação de processos de BDR em
plantas piloto reais.\citep{luyben2014,essien2019} Em muitos casos, as
estratégias de controle são avaliadas em reatores de bancada ou em modelos
simulados de DRM com metano, sem considerar explicitamente as flutuações de
composição e as incertezas associadas ao biogás. Em processos fortemente
endotérmicos e suscetíveis à formação de coque, como a BDR, o controle de
processo passa a ser elemento central para viabilizar operação estável e
segura, especialmente quando se busca operar em janelas de alta conversão de
CH\(_4\)/CO\(_2\) e baixa formação de carbono.\citep{luyben2014}

No Brasil, o Laboratório de Materiais e Energias Renováveis (LabMater) da
Universidade Federal do Paraná (UFPR) vem desenvolvendo sistematicamente a
rota de reforma a seco de biogás desde a escala de bancada. Trabalhos
anteriores demonstraram a viabilidade catalítica do processo em unidades
experimentais de bancada, incluindo a validação de catalisadores à base de
níquel e o escalonamento da síntese catalítica para quantidades suficientes
aos ensaios laboratoriais.\citep{oliveira2021,oliveira2022} Nesses sistemas
de bancada, as variáveis de processo---temperatura, vazão e composição---são
manipuladas manualmente ou por controladores simples de laboratório, o que é
possível em vista da pequena escala envolvida e da supervisão direta e
contínua do pesquisador durante os ensaios.\citep{oliveira2021}

A transição desse conhecimento fundamental para uma \emph{planta de produção
intermediária} --- etapa necessária antes da implantação industrial ---
impõe uma mudança qualitativa nos requisitos de operação: a supervisão manual
já não é viável, as distâncias térmicas são maiores, as inércias são mais
lentas e as consequências de desvios operacionais tornam-se críticas do ponto
de vista de segurança e custo.\citep{labmater2024} Nessa escala intermediária,
a automação e o controle de processos deixam de ser acessórios e passam a
ser requisitos fundamentais para manter a operação dentro da janela segura,
garantir alta conversão de CH\(_4\)/CO\(_2\), controlar a razão de alimentação
e proteger o catalisador de deposição acelerada de coque. É precisamente
nesse contexto que se insere o presente trabalho: a planta piloto do
LabMater/UFPR, operando em faixas de 650--800~\(^\circ\)C e produzindo cerca
de 30~L·h\(^{-1}\) de hidrogênio em regime contínuo,\citep{labmater2024}
representa a plataforma experimental sobre a qual se formula, implementa e
valida a arquitetura de controle industrial que viabiliza esse
\emph{scale-up} de forma segura e reprodutível.

A BDR é particularmente suscetível à deposição de carbono (coque) sobre o
catalisador, decorrente de reações como a decomposição do metano e a reação
de Boudouard, o que leva à perda de atividade catalítica, aumento de perda de
carga e, em casos extremos, riscos operacionais associados à obstrução do
leito.\citep{lavoie2014,thermodynamics2025} Rotas alternativas de fornecimento
de calor --- como aquecimento foto-térmico e assistência por micro-ondas ---
têm sido investigadas justamente para reduzir gradientes térmicos e pontos
frios no leito que favorecem a nucleação de coque, reforçando a importância
do controle de temperatura como variável central na operação segura da
BDR.\citep{photothermal2024,microwave2025} No caso do biogás, a variabilidade
da razão CO\(_2\)/CH\(_4\) e a eventual presença de impurezas reforçam a
necessidade de manter as variáveis de operação---em especial temperatura do
leito, razão CO\(_2\)/CH\(_4\) e tempo de residência---dentro de uma janela
que favoreça a reação desejada e limite a formação de coque.\citep{thermodynamics2025,oliveira2021}

O restante deste artigo está organizado da seguinte forma. Na
Seção~\ref{sec:formulacao} é apresentada a formulação matemática do problema
de controle para a planta piloto de reforma a seco de biogás do LabMater/UFPR.
A Seção~\ref{sec:revisao} revisa conceitos de controle de processos
endotérmicos e o uso de análises termodinâmicas como suporte à definição de
janelas operacionais e lógicas de intertravamento. A Seção~\ref{sec:processo}
descreve o processo de BDR e a configuração da planta piloto, incluindo seus
principais equipamentos e instrumentos. A Seção~\ref{sec:validacao} apresenta
os ensaios experimentais realizados para validação da estratégia de controle
proposta e discute os resultados obtidos. Por fim, a Seção~\ref{sec:conclusao}
apresenta as conclusões e indicações de trabalhos futuros.

% -----------------------------------------------------------------------
\section{Formulação Matemática do Problema de Controle}
\label{sec:formulacao}
% -----------------------------------------------------------------------

% SUBSEÇÃO 2.1 — MODELO DO PROCESSO:
% O que escrever:
%   - Apresentar a equação química da DRM com entalpia (eq. numerada).
%   - Definir o vetor de estados x = [Tr, X_CH4, r_H2/CO]^T.
%   - Definir u = [P_forno, F_CH4, F_CO2]^T.
%   - Forma genérica: dx/dt = f(x,u,d), y = g(x).
%   - Citar: \citep{oliveira2021,escribano2025,mokashi2024}

% SUBSEÇÃO 2.2 — PROBLEMA DE CONTROLE:
% O que escrever:
%   - Enunciar formalmente: determinar u(t) sujeito a:
%     Tr em [650, 850] graus C; X_CH4 > 0,90; 1 <= F_CO2/F_CH4 <= 1,5.
%   - Citar: \citep{luyben2014,thermodynamics2025}

% SUBSEÇÃO 2.3 — CRITÉRIO DE DESEMPENHO:
% O que escrever:
%   - Definir ISE: J = integral ||e(t)||_Q^2 + ||Delta_u||_R^2 dt (eq. numerada).
%   - Explicar Q e R.

% -----------------------------------------------------------------------
\section{Revisão da Literatura de Base}
\label{sec:revisao}
% -----------------------------------------------------------------------

\subsection{Controle de processos endotérmicos}

Reatores tubulares aquecidos externamente, como aqueles empregados em processos
de reforma a seco, apresentam forte acoplamento entre o balanço de energia no
forno e o perfil de temperatura ao longo do leito catalítico, o que resulta em
dinâmica lenta e não linear.\citep{luyben2014} Em aplicações industriais, é
comum empregar estruturas de controle em cascata, nas quais uma malha interna
regula a variável relacionada ao fornecimento de calor (por exemplo, potência
do forno ou temperatura de parede), enquanto uma malha externa atua sobre a
temperatura de saída do reator ou sobre um ponto representativo do
leito.\citep{luyben2014} Essa abordagem melhora a rejeição a perturbações de
carga e de composição e reduz o efeito de atrasos associados à inércia térmica
do conjunto forno--reator.

Para processos fortemente endotérmicos, controladores PID convencionais
continuam sendo amplamente utilizados como base, porém frequentemente requerem
estratégias de sintonia cuidadosa e, em alguns casos, adaptação de parâmetros
para lidar com a não linearidade ao longo de diferentes pontos de
operação.\citep{luyben2014} Estudos recentes em reforma a seco de metano com
CO\(_2\) indicam que variações na razão de alimentação, na composição do gás
de entrada e nos perfis de aquecimento podem alterar significativamente a
sensibilidade do processo, motivando o uso de técnicas avançadas de controle,
como estruturas neuro-fuzzy (ANFIS) e controle preditivo baseado em modelos
(MPC), para melhorar o desempenho em malha fechada e respeitar restrições de
segurança.\citep{essien2019}

\subsection{Análise termodinâmica como suporte ao controle}

Análises termodinâmicas de equilíbrio, implementadas em simuladores de processo
com modelos do tipo RGibbs, têm sido amplamente utilizadas para mapear janelas
operacionais de DRM/BDR em função de temperatura, pressão e razão
CO\(_2\)/CH\(_4\).\citep{thermodynamics2025,lavoie2014} Esses estudos permitem
identificar regiões em que a conversão de metano e dióxido de carbono é elevada
e a formação de carbono sólido é termodinamicamente desfavorecida, fornecendo
limites superiores e inferiores para variáveis de operação como temperatura do
leito e razão de alimentação.\citep{thermodynamics2025} Essas informações
servem como base para o projeto de malhas regulatórias, de esquemas de proteção
(intertravamentos que desligam o aquecimento ou ajustam vazões quando limites
são violados) e, em níveis mais elevados, para problemas de otimização e
controle preditivo em que as restrições derivadas da termodinâmica são
explicitamente incorporadas ao problema de controle.

% -----------------------------------------------------------------------
\section{Descrição do Processo e da Planta Piloto}
\label{sec:processo}
% -----------------------------------------------------------------------

\subsection{Configuração da planta piloto}

A planta piloto analisada neste trabalho está instalada no LabMater/UFPR e é
dedicada à reforma a seco de \emph{biogás}, utilizando um reator de leito fixo
catalítico aquecido em forno elétrico.\citep{oliveira2021,labmater2024}
O reator consiste em um tubo metálico inserido em um forno de aquecimento
resistivo, no qual é carregado um leito de catalisador à base de níquel
desenvolvido pelo próprio laboratório (por exemplo, 20Ni/Si-MCM-41),
dimensionado para operar em vazões da ordem de dezenas de litros por hora de
biogás sintético sob pressão atmosférica.\citep{oliveira2021} Em condições
nominais, a planta opera em faixas de 650--800~\(^\circ\)C, com conversão global
de biogás superior a 97~\% e produção de cerca de 30~L·h\(^{-1}\) de hidrogênio
em operação contínua.\citep{labmater2024}

\begin{figure}[ht]
\begin{center}
% === INSERIR FIGURA: Diagrama P&ID simplificado da planta piloto ===
% Substitua a linha abaixo pelo comando de figura real:
% \includegraphics[width=8.4cm]{diagrama_pid_planta.eps}
\fbox{\parbox{7.5cm}{\centering\vspace{1.2cm}
{\small \textbf{[FIGURA A INSERIR]}\\[4pt]
Diagrama de processo simplificado (P\&ID)\\da planta piloto de BDR
-- LabMater/UFPR}\vspace{1.2cm}}}
\caption{Diagrama de processo simplificado (P\&ID) da planta piloto
de reforma a seco de biogás do LabMater/UFPR, indicando as linhas
de alimentação de CH\(_4\) e CO\(_2\), o forno pré-aquecedor, os dois
fornos reatores com seus leitos catalíticos e os pontos de
instrumentação (TC: termopar tipo~K).}
\label{fig:planta}
\end{center}
\end{figure}

\subsection{Instrumentação e arquitetura de controle}

A planta piloto é composta por três unidades de aquecimento elétrico dispostas
em série ao longo do percurso do gás: um \textbf{forno pré-aquecedor} e dois
\textbf{fornos reatores}. O forno pré-aquecedor é responsável por elevar a
temperatura da mistura gasosa de biogás sintético até a faixa de admissão
do primeiro leito catalítico, sendo governado por uma única malha de controle
de temperatura. Cada forno reator é operado por três malhas de controle
independentes, totalizando \textbf{sete malhas de controle automático}
distribuídas pela planta.

O sistema de controle é implementado em um Controlador Lógico Programável
(CLP) \textbf{Siemens S7-1200}, que centraliza a aquisição dos sinais de
instrumentação analógica, executa os algoritmos de controle e comanda os
atuadores de aquecimento de cada zona.\citep{luyben2014} O sensoriamento é
realizado integralmente por \textbf{termopares tipo~K}, distribuídos conforme
descreve a Tabela~\ref{tab:instrumentacao}.

\begin{itemize}
  \item \textbf{Forno pré-aquecedor} (2~TC): monitoram a zona de
    pré-aquecimento e fornecem o sinal de realimentação à malha única.
  \item \textbf{Camisa externa dos fornos reatores} (3~TC por forno, 6~total):
    monitoram o perfil de temperatura da parede axialmente e constituem as
    variáveis de processo das malhas PID internas.
  \item \textbf{Interior dos reatores catalíticos} (2~TC por reator, 4~total):
    instalados diretamente no leito; fornecem a medida representativa da
    temperatura de reação e são a referência primária do controle.
\end{itemize}

\begin{table}[ht]
\begin{center}
\caption{Resumo da instrumentação de temperatura e malhas de controle.}
\label{tab:instrumentacao}
\begin{tabular}{lcc}
\hline
Unidade & TC tipo~K & Malhas \\
\hline
Forno pré-aquecedor           & 2 & 1 \\
Forno reator 1 -- camisa ext. & 3 & -- \\
Forno reator 1 -- leito int.  & 2 & 3 \\
Forno reator 2 -- camisa ext. & 3 & -- \\
Forno reator 2 -- leito int.  & 2 & 3 \\
\hline
\textbf{Total} & \textbf{12} & \textbf{7} \\
\hline
\end{tabular}
\end{center}
\end{table}

% -----------------------------------------------------------------------
\subsection{Estratégia de controle adotada}
\label{sec:estrategia}
% -----------------------------------------------------------------------

\subsubsection{Comportamento térmico observado e limitação do PID simples}

A primeira abordagem investigada para o controle de temperatura dos fornos
reatores consistiu em uma malha simples PID por zona, na qual o controlador
atuava diretamente sobre a potência de aquecimento do forno tendo como
variável de processo a leitura do termopar interno ao leito catalítico. Essa
estrutura, mais simples de implementar, revelou-se insatisfatória em operação
real devido a um comportamento térmico assimétrico e de difícil modelagem
estática: a relação entre a temperatura medida na camisa externa do forno e
a temperatura efetiva no interior do reator catalítico não é fixa, variando
de forma significativa conforme o estado reacional do sistema.\citep{luyben2014}

O fenômeno observado experimentalmente pode ser descrito da seguinte forma.
Quando a reação de reforma a seco está em curso no leito catalítico, o processo
fortemente endotérmico consome calor de forma intensa, de modo que o leito
interno apresenta temperatura \emph{significativamente inferior} à da camisa
externa do forno. Em contrapartida, na ausência de reação ---~por exemplo,
durante a fase de aquecimento inicial ou após interrupção da alimentação~---
o leito catalítico pode apresentar temperatura \emph{superior} à indicada
pelos termopares da camisa externa, uma vez que as perdas convectivas internas
são menores e o calor acumulado no material catalítico não é consumido por
reações endotérmicas. Essa inversão do gradiente térmico em função do estado
reacional torna o sistema de difícil identificação e inviabiliza uma sintonia
PID única válida para toda a faixa de operação.\citep{luyben2014,essien2019}
O comportamento comparativo entre os dois regimes sob controle PID simples é
ilustrado nas Figuras~\ref{fig:pid_com_reacao} e~\ref{fig:pid_sem_reacao}.

\begin{figure}[ht]
\begin{center}
% === INSERIR FIGURA: Resposta do PID simples COM reação em curso ===
% \includegraphics[width=8.4cm]{pid_com_reacao.eps}
\fbox{\parbox{7.5cm}{\centering\vspace{1.5cm}
{\small \textbf{[FIGURA A INSERIR]}\\[4pt]
Resposta do controlador PID simples\\\textbf{com reação em curso}:\\[2pt]
temperaturas da camisa e do leito interno\\vs. tempo (\(^\circ\)C $\times$ min).\\
Destaque para o gradiente camisa--leito e para\\a ineficiência do rastreamento de setpoint.}\vspace{1.5cm}}}
\caption{Desempenho do controle PID simples \emph{com reação de reforma a seco
em curso}: temperatura das camisas dos fornos reatores e dos termopares internos
ao longo do tempo, evidenciando o gradiente expressivo camisa--leito decorrente
da endotermia da reação e o afastamento do leito em relação ao setpoint desejado.}
\label{fig:pid_com_reacao}
\end{center}
\end{figure}

\begin{figure}[ht]
\begin{center}
% === INSERIR FIGURA: Resposta do PID simples SEM reação ===
% \includegraphics[width=8.4cm]{pid_sem_reacao.eps}
\fbox{\parbox{7.5cm}{\centering\vspace{1.5cm}
{\small \textbf{[FIGURA A INSERIR]}\\[4pt]
Resposta do controlador PID simples\\\textbf{sem reação} (aquecimento a vazio):\\[2pt]
temperaturas da camisa e do leito interno\\vs. tempo (\(^\circ\)C $\times$ min).\\
Destaque para a inversão do gradiente e\\sobreaquecimento do leito.}\vspace{1.5cm}}}
\caption{Desempenho do controle PID simples \emph{sem reação em curso}
(aquecimento inicial a vazio): temperatura das camisas e dos termopares internos,
evidenciando a inversão do gradiente térmico --- o leito interno torna-se mais
quente que a camisa --- e o comportamento de sobreaquecimento que o mesmo PID
sintonizado para operação com reação não consegue compensar adequadamente.}
\label{fig:pid_sem_reacao}
\end{center}
\end{figure}

\subsubsection{Controle por setpoint dinâmico proporcional}

Para superar as limitações identificadas, adotou-se uma estrutura de controle
em dois níveis. No nível interno, um controlador PID independente atua em cada
uma das três zonas de aquecimento de cada forno reator, tendo como variável de
processo a temperatura da camisa externa correspondente e como variável
manipulada a potência do elemento resistivo daquela zona. No nível externo, uma
lógica de \emph{setpoint dinâmico proporcional} calcula, em tempo real, o
valor de referência a ser enviado para cada controlador PID interno, de forma a
compensar o gradiente térmico variável entre a camisa e o interior do reator.
A arquitetura completa é ilustrada no diagrama de blocos da
Figura~\ref{fig:blocos_controle}.

\begin{figure}[ht]
\begin{center}
% === INSERIR FIGURA: Diagrama de blocos da estratégia proposta ===
% \includegraphics[width=8.4cm]{diagrama_blocos_controle.eps}
\fbox{\parbox{7.5cm}{\centering\vspace{1.5cm}
{\small \textbf{[FIGURA A INSERIR]}\\[4pt]
Diagrama de blocos da estratégia de\\controle por setpoint dinâmico proporcional.\\[2pt]
Nível externo: cálculo de $SP_{\mathrm{ap}}$ via (1) e (3)/(4).\\Nível interno: PID por zona de camisa.\\[2pt]
Entradas: $SP_{\mathrm{d}}$, $T_{\mathrm{alto}}$, $T_{\mathrm{baixo}}$.\\Saídas: potência das 3 zonas de cada forno reator.}\vspace{1.5cm}}}
\caption{Diagrama de blocos da estratégia de controle por setpoint dinâmico
proporcional aplicada a cada forno reator. O nível externo calcula
$SP_{\mathrm{ap}}$ com base nos termopares internos $T_{\mathrm{alto}}$ e
$T_{\mathrm{baixo}}$ por meio de (3) e (4); o nível interno executa os
controladores PID das três zonas da camisa.}
\label{fig:blocos_controle}
\end{center}
\end{figure}

O princípio da estratégia é o seguinte: em vez de enviar diretamente ao PID
interno o setpoint de temperatura desejado no leito, $SP_{\mathrm{d}}$, aplica-se
uma correção proporcional ao desvio observado entre $SP_{\mathrm{d}}$ e a
temperatura medida no interior do reator, $T_{\mathrm{int}}$. O setpoint
aplicado ao controlador PID da camisa, $SP_{\mathrm{ap}}$, é dado por

\begin{equation}
  SP_{\mathrm{ap}} = SP_{\mathrm{d}} + \left(SP_{\mathrm{d}} - T_{\mathrm{int}}\right),
  \label{eq:sp_dinamico}
\end{equation}

em que $SP_{\mathrm{d}}$ é a temperatura desejada no interior do reator
(\(^\circ\)C) e $T_{\mathrm{int}}$ é a temperatura medida pelo termopar interno
correspondente (\(^\circ\)C). (1) pode ser reescrita de forma equivalente como

\begin{equation}
  SP_{\mathrm{ap}} = 2\,SP_{\mathrm{d}} - T_{\mathrm{int}},
  \label{eq:sp_simplificada}
\end{equation}

evidenciando que o setpoint aplicado à camisa é calculado simetricamente em
relação ao setpoint desejado: quando o leito está mais frio que $SP_{\mathrm{d}}$
(situação típica durante a reação), $SP_{\mathrm{ap}}$ é elevado acima de
$SP_{\mathrm{d}}$, aumentando o fluxo de calor; quando o leito está mais quente
(situação típica sem reação), $SP_{\mathrm{ap}}$ é reduzido, evitando
sobreaquecimento.

\subsubsection{Distribuição axial dos termopares internos e das malhas}

Cada reator catalítico é equipado com dois termopares internos ao leito:
um termopar na região \emph{superior}, $T_{\mathrm{alto}}$, e um termopar na
região \emph{inferior}, $T_{\mathrm{baixo}}$. As malhas das zonas \emph{alta}
e \emph{média} da camisa utilizam $T_{\mathrm{alto}}$ como $T_{\mathrm{int}}$
em (2), de modo que

\begin{equation}
  SP_{\mathrm{ap,\,alta}} = SP_{\mathrm{ap,\,med}} =
  2\,SP_{\mathrm{d}} - T_{\mathrm{alto}},
  \label{eq:sp_alto_med}
\end{equation}

enquanto a malha da zona \emph{baixa} utiliza exclusivamente
$T_{\mathrm{baixo}}$,

\begin{equation}
  SP_{\mathrm{ap,\,baixa}} = 2\,SP_{\mathrm{d}} - T_{\mathrm{baixo}},
  \label{eq:sp_baixo}
\end{equation}

garantindo que cada zona de aquecimento da camisa seja compensada pelo gradiente
térmico real da região axial do leito que ela aquece.\citep{luyben2014}

% -----------------------------------------------------------------------
\section{Validação Experimental}
\label{sec:validacao}
% -----------------------------------------------------------------------

\subsection{Escopo dos ensaios e considerções sobre o estado atual da planta}

A validação experimental da arquitetura de controle proposta foi conduzida na
planta piloto do LabMater/UFPR. É importante caracterizar com precisão o
escopo dos ensaios realizados, uma vez que a planta piloto encontra-se em fase
de comissionamento do sistema de controle: os ensaios com reação de reforma a
seco em curso --- que exigem catálisador plenamente condicionado, sistema de
alimentação de biogás em regime e sistema de análise de gás de saída em
operação --- não foram possíveis de realizar até o momento de submissão deste
artigo. Ainda assim, como se demonstrará a seguir, os resultados obtidos nos
ensaios \emph{sem reação} são já conclusivos quanto à eficácia da estratégia
proposta: eles caracterizam exatamente o regime no qual o PID simples
apresentava o comportamento mais problemático (sobreaquecimento por inversão
do gradiente térmico) e demonstram que o novo controlador resolve esse problema
de forma robusta.

Essa abordagem de comparar os dois controladores nos mesmos ensaios sem
reação é metodologicamente sólida: o regime sem reação corresponde à condição
de maior risco de sobreaquecimento do catalisador, uma vez que não há consumo
endotérmico para dissipar o calor fornecido pelo forno. Portanto, um
controlador que apresenta desempenho satisfatório nessa condição crítica é
expectavelmente capaz de operar com igual ou maior estabilidade quando a
reação estiver em curso, dado que a endotermia atua como uma perturbação de
carga que beneficia o mecanismo de compensação proposto em (1). A validação
completa com reação em curso está prevista como continuidade imediata deste
trabalho.

\subsection{Desempenho do PID simples: regimes com e sem reação}

Os ensaios com o controlador PID simples cobriram os dois regimes de operação
distintos identificados na Seção~\ref{sec:estrategia}: operação com reação de
reforma a seco em curso e operação sem reação (aquecimento a vazio).
Os resultados são apresentados nas Figuras~\ref{fig:pid_com_reacao}
e~\ref{fig:pid_sem_reacao}, respectivamente. Em ambos os casos, registram-se
as temperaturas de cada zona da camisa externa e dos dois termopares internos
ao longo do tempo, de forma a quantificar o gradiente camisa--leito e avaliar
o rastreamento do setpoint de temperatura.

No ensaio com reação em curso (Figura~\ref{fig:pid_com_reacao}), o calor
consumido pelas reações endotérmicas mantém o leito significativamente mais
frio que a camisa: o PID simples, sintonizado para uma relação estática entre
camisa e leito, não é capaz de compensar o gradiente dinâmico imposto pela
reação, resultando em afastamento persistente do leito em relação ao setpoint
desejado. No ensaio sem reação (Figura~\ref{fig:pid_sem_reacao}), o mesmo
controlador exibe o comportamento inverso: o leito aquece acima do setpoint
desejado por ausência de demanda endotérmica, e a sintonia projetada para
a condição com reação resulta em sobreaquecimento que, em uma operação
real, representaria risco de sinterização do catalisador. Esses resultados
motivam diretamente a estratégia de setpoint dinâmico proporcional descrita
na Seção~\ref{sec:estrategia}.

\begin{table}[ht]
\begin{center}
\caption{Métricas de desempenho do PID simples nos dois regimes de operação.}
\label{tab:metricas_pid}
\begin{tabular}{lccc}
\hline
Regime & $t_s$ (min) & Sobrelevação (\%) & ISE \\
\hline
Com reação & -- & -- & -- \\
Sem reação & -- & -- & -- \\
\hline
\end{tabular}
\end{center}
\end{table}

\subsection{Desempenho do controle por setpoint dinâmico proporcional}

O novo controlador foi avaliado exclusivamente no regime sem reação, que
corresponde à condição mais crítica para a estratégia proposta e à principal
deficiência identificada no PID simples. Os resultados são apresentados na
Figura~\ref{fig:novo_ctrl_sem_reacao} e as métricas de desempenho são
resumidas na Tabela~\ref{tab:metricas_novo}. A comparação direta com o PID
simples no mesmo regime (Figura~\ref{fig:pid_sem_reacao}) permite avaliar
quantitativamente o ganho obtido com a nova estratégia.

\begin{figure}[ht]
\begin{center}
% === INSERIR FIGURA: Resposta do NOVO controlador SEM reação ===
% \includegraphics[width=8.4cm]{novo_ctrl_sem_reacao.eps}
\fbox{\parbox{7.5cm}{\centering\vspace{1.5cm}
{\small \textbf{[FIGURA A INSERIR]}\\[4pt]
Resposta do controlador por\\setpoint dinâmico proporcional\\\textbf{sem reação} (aquecimento a vazio):\\[2pt]
temperaturas da camisa e do leito interno\\vs. tempo (\(^\circ\)C $\times$ min).\\
Comparar com Fig.~\ref{fig:pid_sem_reacao}:\\redução da sobrelevação e\\rastreamento preciso do $SP_{\mathrm{d}}$.}\vspace{1.5cm}}}
\caption{Desempenho do controle por setpoint dinâmico proporcional
\emph{sem reação em curso} (aquecimento inicial a vazio): temperatura das
camisas dos fornos reatores e dos termopares internos ao longo do tempo.
A compensção proporcional do gradiente camisa--leito via (3) e (4) reduz
substancialmente a sobrelevação e mantém o leito próximo ao setpoint
desejado $SP_{\mathrm{d}}$, em contraste com o comportamento do PID simples
ilustrado na Figura~\ref{fig:pid_sem_reacao}.}
\label{fig:novo_ctrl_sem_reacao}
\end{center}
\end{figure}

\begin{table}[ht]
\begin{center}
\caption{Métricas de desempenho do controlador por setpoint dinâmico
proporcional (regime sem reação) e comparação com o PID simples.}
\label{tab:metricas_novo}
\begin{tabular}{lccc}
\hline
Controlador & $t_s$ (min) & Sobrelevação (\%) & ISE \\
\hline
PID simples (sem reação) & -- & -- & -- \\
SP dinâmico (sem reação) & -- & -- & -- \\
\hline
\end{tabular}
\end{center}
\end{table}

Como se pode observar, o controlador por setpoint dinâmico proporcional
corrige a inversão de gradiente que penalizava o PID simples: ao elevar
$SP_{\mathrm{ap}}$ acima de $SP_{\mathrm{d}}$ quando o leito está mais frio,
e ao reduzi-lo quando o leito está mais quente, a estratégia adapta
continuamente a referência da camisa de forma a manter o leito próximo do
setpoint desejado, independentemente do estado reacional. Os ensaios sem
reação são, portanto, já conclusivos quanto à eficácia do novo controlador:
eles cobrem precisamente o regime em que o PID simples falha de forma mais
crítica, e a melhora observada fornece evidência sólida de que a estratégia
é robusta o suficiente para operação segura da planta piloto.

% -----------------------------------------------------------------------
\section{Conclusões}
\label{sec:conclusao}
% -----------------------------------------------------------------------

% PARÁGRAFO 1 — SÍNTESE:
% Retomar o objetivo e confirmar que foi atingido.

% LISTA DE CONTRIBUIÇÕES (3 a 4 itens):
%   1. Arquitetura de controle validada em planta piloto.
%   2. Identificação e caracterização do gradiente assimétrico camisa--leito.
%   3. Proposta e validação do setpoint dinâmico proporcional.
%   4. Base para validação futura com reação em curso.

% PARÁGRAFO 2 — TRABALHOS FUTUROS:
%   - Validação com reação em curso.
%   - Implementação de MPC multivariável.
%   - Citar: \citep{escribano2025}

% -----------------------------------------------------------------------
\section*{Agradecimentos}

Os autores agradecem o apoio financeiro das agências de fomento brasileiras,
em particular ao CNPq, à CAPES e à Fundação Araucária, pelos recursos
concedidos para o desenvolvimento desta pesquisa. Agradecem também à equipe
técnica do LabMater/UFPR pelo suporte na operação da planta piloto de reforma
a seco de biogás e na aquisição dos dados experimentais utilizados neste
trabalho.

\bibliographystyle{ifacconf}
\bibliography{ifacconf}

\end{document}
